\labsection{9}{Validate an imbalanced classifier without leakage}{sec:lab09}

\begin{labproject}{Validation and preprocessing: evaluate the process, not one lucky split}
\textbf{Coverage.} Chapter~\ref{ch:validation}.\\
\textbf{Goal.} You should be able to construct stratified folds, prevent preprocessing leakage, compare thresholds across folds, quantify metric variability, and distinguish discrimination from calibration.\\
\textbf{Datasets.} \texttt{credit\_card\_fraud\_dataset.csv}; Heart Failure is used for a calibration extension.\\
\textbf{Scenario.} Fraud is rare, so a superficially high accuracy can be operationally useless. You must evaluate a complete modelling process with stratified cross-validation, fold-local preprocessing, and threshold-sensitive metrics.

\textbf{Core workflow.}
\begin{enumerate}
  \item \textbf{Measure the class imbalance.} Report fraud prevalence and the accuracy of an ``always legitimate'' classifier. Treat this as the minimum baseline any useful model must beat.
  \item \textbf{Create stratified folds.} Build five folds that each contain fraud and legitimate transactions. Report the class count in every fold.
  \item \textbf{Keep preprocessing inside the fold.} For each validation fold, estimate scaling parameters from the corresponding training folds only, then apply those parameters to validation data.
  \item \textbf{Fit and retain probabilities.} Train logistic regression inside each fold and store validation probabilities so the same predictions can be evaluated at multiple thresholds.
  \item \textbf{Evaluate decision thresholds.} Compare 0.10, 0.25, and 0.50 using accuracy, precision, recall, and specificity. Include false-positive and false-negative counts where useful.
  \item \textbf{Report variability.} Summarise both the mean and fold-to-fold spread of the metrics. Identify whether recall is stable enough to support a reliable operational claim.
  \item \textbf{Choose a threshold from a stated cost.} Define a hypothetical relative cost for missed fraud versus false alarms, then select the threshold that best reflects that priority. Make clear that the cost assumption is external to the model.
\end{enumerate}

\textbf{Calibration extension.} Reuse the Heart Failure modelling workflow, bin held-out predicted probabilities, and compare mean predicted probability with observed event rate. Explain why a model can rank cases well yet still produce poorly calibrated probabilities.

\tcbbreak
\lstinputlisting[style=utpr]{all_experiments/lab09_validation.R}
\begin{verification}
\textbf{Required evidence.} Submit the majority-class baseline, five-fold metric summary, one variability result, and a threshold recommendation tied to an explicit error-cost assumption. Include one sentence identifying where leakage would occur if global scaling were applied before cross-validation.

\textbf{Professional standard.} Validation must reproduce the full future prediction process. Any transformation learned from all records before validation contaminates the estimate of generalisation performance.
\end{verification}
\end{labproject}
